% Part: first-order-logic
% Chapter: natural-deduction
% Section: rules-and-proofs

\documentclass[../../../include/open-logic-section]{subfiles}

\begin{document}

\iftag{FOL}
      {\olfileid[fr]{fol}{ntd}{rul}}
      {\olfileid[fr]{pl}{ntd}{rul}}

\olsection{Règles et \usetoken{p}{derivation}}

\begin{explain}
Les systèmes de déduction naturelle visent à reproduire au plus près
le raisonnement informel des démonstrations mathématiques, d'où leur
caractère « naturel ». Les preuves en déduction naturelle commencent
par des hypothèses, auxquelles on applique ensuite des règles
d'inférence. Les règles \Intro{\lnot}, \Intro{\lif},
\iftag{FOL}{\Elim{\lor}, \Elim{\lexists}}{\Elim{\lor}} et \FalseCl{}
permettent de « !!{discharge} » des hypothèses. Pour indiquer
clairement quelle inférence décharge une hypothèse, on place à côté
de l'inférence l'étiquette de cette hypothèse.
\end{explain}

\begin{defn}[Hypothèse]
Une \emph{hypothèse} est !!a{sentence} qui occupe la position
la plus haute d'une branche de l'arbre.
\end{defn}

Les !!{derivation}s en déduction naturelle sont certains arbres
d'!!{sentence}s. Les !!{sentence}s situés tout en haut sont les
hypothèses; si !!a{sentence} se trouve sous un, deux ou trois autres
!!{sentence}s, il doit en résulter par une application correcte d'une
règle d'inférence. Les !!{sentence}s au-dessus du trait d'inférence
sont les \emph{prémisses}, et l'!!{sentence} situé en dessous est la
\emph{conclusion} de l'inférence. Les règles vont par paires : une
règle d'introduction et une règle d'élimination pour chaque
!!{operator}. Elles introduisent un !!{operator} dans la conclusion
ou retirent un !!{operator} d'une prémisse. Certaines règles
permettent de \emph{!!{discharge}} une hypothèse d'un type donné.
Pour indiquer quelle inférence décharge quelle hypothèse, on attribue
une même étiquette à l'hypothèse et à l'inférence. On écrit alors
l'hypothèse sous la forme «~$\Discharge{!A}{n}$~».

% N'inclure cette phrase que si l'un des symboles land, lor, lif ou lnot est défini.

\iftag{notprvNot,notprvAnd,notprvOr,notprvIf}{}%
    {On considère habituellement des règles pour tous les !!{operator}s $\land$, $\lor$, $\lif$, $\lnot$ et $\lfalse$, même si certains d'entre eux sont définis à partir des autres.}

\Needspace{7\baselineskip}
\begin{editorial}
La liste initiale de la source omet \FalseCl, qui permet pourtant
de décharger une hypothèse, comme le montre sa règle ci-dessous.
Dans la description des arbres, « séquents » a été corrigé en
« énoncés », conformément aux objets de ce système. L'organisation
en introduction et élimination décrit les connecteurs usuels;
les deux règles particulières pour $\lfalse$ sont présentées
séparément dans la section suivante.
\end{editorial}

\end{document}
